\section{Metodologia}
\label{sec:metodologia}

O trabalho segue a metodologia de \cite{courb}, adaptada aos dados
disponíveis em Belo Horizonte, e acrescenta uma análise do atraso em relação ao horário
programado. O processo tem cinco etapas: coleta, pré-processamento, cálculo dos horários de
passagem nas paradas, réplica dos modelos e extensão. Todo o código está disponível no
repositório do trabalho\footnote{\url{https://github.com/JVDS17/tcc}} e reproduz, a partir dos
dados brutos, todas as tabelas e figuras deste texto.

\subsection{Dados}

\textbf{Posições dos veículos.} As posições de GPS foram coletadas por um programa em Python
que consultou, a cada 20 segundos, o \textit{feed} GTFS-Realtime de posições de veículos
(\textit{VehiclePositions}) do sistema de transporte de Belo Horizonte, entre 30/03/2026 e
19/04/2026. O coletor era executado em sessões de 12 horas e gravava os registros das linhas
monitoradas em uma planilha do Google Sheets. Cada registro contém o instante da consulta, a
linha, o veículo, o identificador da viagem no GTFS (\textit{trip\_id}), o horário de partida
programado, a latitude e a longitude; o campo de velocidade veio sempre zerado e não foi
usado. Foram coletados 105.207 registros da linha 4110, em 16 dias. Houve dias sem coleta (01
a 05/04) e períodos com poucos registros, devidos a instabilidades do \textit{feed}, a falhas
de transmissão dos veículos e aos intervalos entre as sessões de coleta.

\textbf{Quadro de horários e traçado.} Usou-se o GTFS estático da BHTRANS \cite{gtfsbh} na
versão baixada pelo próprio coletor em 19/04/2026 (versão 202604181953), em vigor durante a
coleta. Os 95 \textit{trip\_ids} encontrados na coleta existem nesse GTFS. A escolha da versão
importa: versões posteriores mantêm os mesmos identificadores, mas com horários programados
diferentes em cerca de um terço das paradas. Do GTFS foram extraídos o traçado da linha
(\textit{shapes.txt}), as paradas (\textit{stops.txt}) e os horários programados em cada
parada (\textit{stop\_times.txt}).

\textbf{Recorte.} Todas as viagens registradas pertencem ao sentido Belvedere $\rightarrow$
Dom Cabral, que tem 64 paradas e 19,5~km; o sentido contrário não foi capturado pela coleta.
A análise se restringe, portanto, a esse sentido.

\subsection{Pré-processamento}

\textbf{Viagem real.} O \textit{trip\_id} da BHTRANS identifica um horário de partida do
quadro (por exemplo, a partida das 12h52 em dias úteis) e se repete todos os dias. Por isso,
uma viagem realizada foi definida pela combinação de data, \textit{trip\_id} e veículo, o que
resultou em 995 viagens.

\textbf{Posições repetidas.} Embora a API fosse consultada a cada 20~s, os veículos atualizam
a posição com intervalo mediano de 46~s. Leituras consecutivas do mesmo veículo com a mesma
posição foram tratadas como duplicatas e removidas, restando 41.643 registros.

\subsection{Horário de passagem nas paradas}

O GPS registra posições em instantes discretos, que raramente coincidem com a passagem pela
parada. Para estimar esse instante, cada posição foi projetada sobre o traçado da linha,
obtendo-se a distância percorrida $d$ desde o início do percurso. A busca da projeção foi
limitada a uma janela à frente da posição anterior, evitando associar o ônibus a um trecho
errado em locais onde a rota passa perto de si mesma, e posições a mais de 80~m do traçado
foram descartadas. A distância de cada parada $s$ ao longo da rota, $d_s$, vem do próprio GTFS.

O horário de passagem foi obtido por interpolação linear entre as duas posições consecutivas
$P_i$ e $P_{i+1}$ que envolvem a parada ($d_i \le d_s \le d_{i+1}$):
\begin{equation}
T_s = T_i + (T_{i+1} - T_i)\,\frac{d_s - d_i}{d_{i+1} - d_i}
\end{equation}
Interpolações em que $T_{i+1} - T_i > 180$~s foram descartadas, por indicarem perda de sinal;
esse limite cobre 99\% dos intervalos reais de atualização. Para a primeira parada, usou-se o
instante em que o ônibus se afasta 30~m do terminal, isto é, a partida efetiva.

A API associa \textit{trip\_ids} a muitos fragmentos curtos de trajeto (veículos que aparecem
por poucos minutos). Por isso, mantiveram-se apenas as viagens com horário de passagem
calculado em pelo menos 50\% das paradas. Quando dois veículos apareceram com o mesmo
\textit{trip\_id} no mesmo dia, manteve-se o de maior cobertura. Também foram corrigidas as
viagens rotuladas com o dia de serviço errado na virada da meia-noite. Ao final, restaram 393
viagens e 21.871 horários de passagem, com cobertura média de 87\% das paradas.

\subsection{Réplica: tempo de deslocamento entre paradas}

Seguindo \cite{courb}, a variável resposta é o tempo de deslocamento
$t_{jk}$ no trecho $j$ (da parada $j$ à parada $j+1$) na viagem $k$. Os trechos com menos de
100 observações (a saída do terminal e a chegada final) foram excluídos, restando 60 trechos e
21.372 observações. Foram ajustados quatro modelos:

\begin{itemize}
  \item \textbf{M1}, apenas efeito fixo: $t_{jk} = \beta_0 + \varepsilon_{jk}$;
  \item \textbf{M2}, com intercepto aleatório por trecho:
        $t_{jk} = \beta_0 + b_{j} + \varepsilon_{jk}$, com $b_j \sim N(0, \sigma_b^2)$ e
        $\varepsilon_{jk} \sim N(0, \sigma^2)$;
  \item \textbf{M3}, igual ao M2, mas com erro autorregressivo de primeira ordem dentro de cada
        trecho: $\varepsilon_{jk} = \phi\,\varepsilon_{j,k-1} + \eta_{jk}$, em que $k-1$ é a
        viagem anterior que passou pelo mesmo trecho;
  \item \textbf{M3 sazonal}, em que os resíduos do M3, ordenados cronologicamente, recebem um
        modelo AR(1) com componente sazonal AR(1) de período igual ao número de trechos por
        viagem (60).
\end{itemize}

Os modelos M2 e M3 foram ajustados por máxima verossimilhança com o pacote \texttt{nlme} do R
\cite{pinheiro2000}, e o componente sazonal com o pacote \texttt{forecast}
\cite{hyndman2008}. O código partiu do repositório público de \cite{courb}, com três diferenças: o coeficiente $\phi$ é estimado a partir dos dados, e
não fixado; o componente sazonal é estimado por soma de quadrados condicional, porque a
máxima verossimilhança exata não convergiu em uma série com cerca de 24 mil pontos, período 60
e lacunas; e todas as viagens da linha entram no modelo, e não apenas as de um veículo.

\textbf{Avaliação.} Como no artigo original, compara-se o erro quadrático médio (EQM) no
ajuste. Além disso, os modelos foram avaliados fora da amostra de duas formas: (a) ajuste com
os dados de 30/03 a 12/04 e teste com os de 13 a 19/04; e (b) ajuste com a quarta-feira 08/04
e teste com a quarta-feira 15/04, reproduzindo o desenho do artigo. Em ambas, as previsões do
M3 e do M3 sazonal usam os resíduos já observados (previsão um passo à frente), e são
reportados a raiz do erro quadrático médio (RMSE) e o erro absoluto médio (MAE).

\subsection{Extensão: atraso em relação ao horário programado}

O atraso na parada $j$ da viagem $i$ é a diferença entre o horário de passagem estimado e o
programado no GTFS:
\begin{equation}
A_{ij} = T_{ij} - T^{prog}_{ij}
\end{equation}
Valores positivos indicam atraso. No GTFS da BHTRANS, apenas a primeira parada é ponto de
controle; os horários das demais são estimativas da própria operadora. O atraso nas paradas
intermediárias é, portanto, medido em relação a essa estimativa.

Foram ajustados três modelos lineares mistos com o pacote \texttt{lme4} \cite{bates2015}, por
máxima verossimilhança restrita (REML):

\begin{itemize}
  \item \textbf{E1}: $A_{ij} = \beta_0 + \beta_1\,\mathrm{pico}_i + u_i + \varepsilon_{ij}$,
        com $u_i$ o efeito aleatório da viagem real e $\mathrm{pico}_i = 1$ para partidas
        programadas entre 16h e 19h59 em dias úteis. É a formulação proposta na primeira etapa
        deste trabalho.
  \item \textbf{E1b}: igual ao E1, mas com o efeito aleatório pelo \textit{trip\_id}. Foi
        incluído para mostrar o efeito de agrupar a mesma partida de dias diferentes, como
        havia sido feito na primeira etapa.
  \item \textbf{E2}: $A_{ij} = \beta_0 + \boldsymbol{\beta}_f\,\mathrm{faixa}_i +
        \boldsymbol{\beta}_d\,\mathrm{tipo\_dia}_i + \beta_k\,\mathrm{km}_{j} + u_i + v_j +
        \varepsilon_{ij}$, com faixa horária da partida (até 9h, 9h--16h, 16h--20h e após
        20h), tipo de dia (útil, sábado e domingo), distância da parada ao início do percurso
        (km), efeito aleatório da viagem $u_i$ e da parada $v_j$.
\end{itemize}

Como o número de viagens completas no pico da tarde foi pequeno, os modelos E1 e E2 foram
reajustados com um critério de cobertura menos rígido (25\% das paradas) para verificar a
sensibilidade do efeito do pico a esse critério.
